% Partie : intuitionistic-logic
% Chapitre : tableaux
% Section : soundness

\documentclass[../../../include/open-logic-section]{subfiles}

\begin{document}

\olfileid[fr]{int}{tab}{sou}

\olsection{Correction des \usetoken{p}{tableau} intuitionnistes}

\begin{explain}
  Pour montrer que les !!{tableau}s intuitionnistes sont corrects, il
  faut établir que, si les hypothèses
  \[
  \sFmla{\True}{!B_1}[1], \dots, \sFmla{\True}{!B_n}[1], \sFmla{\False}{!A}[1]
  \]
  admettent un !!{tableau} fermé, alors $!B_1, \dots, !B_n \Entails
  !A$. Il est plus simple de prouver la contraposée : s'il existe un
  modèle~$\mModel{M}$ et un monde~$w$ tels que
  $\mSat{M}{!B_i}[w]$ pour tout $i=1$, \dots,~$n$, mais
  $\mSat/{M}{!A}[w]$, aucun !!{tableau} issu de ces hypothèses ne
  peut être fermé. Ce contre-modèle montre que les hypothèses
  initiales du !!{tableau} sont satisfaisables. La stratégie
  consiste à montrer que, si toutes les !!{formula}s préfixées d'une
  branche d'!!a{tableau} sont satisfaisables, l'application d'une
  règle laisse au moins une branche prolongée satisfaisable. Puisque
  les branches fermées ne sont pas satisfaisables, tout !!{tableau}
  construit à partir d'!!{formula}s préfixées satisfaisables possède
  au moins une branche ouverte.

  Pour appliquer cette stratégie aux tableaux intuitionnistes, nous
  devons définir la satisfaisabilité des formules signées munies de
  préfixes relativement à un modèle relationnel. La preuve découle
  ensuite des règles et de la persistance de la vérité.
\end{explain}

\begin{defn}
  Soit $P$ un ensemble de préfixes, c'est-à-dire $P \subseteq
  (\PosInt)^* \setminus \{\emptyseq\}$, et soit $\mModel{M}$ un
  modèle. Une fonction~$f\colon P \to W$ est une
  \emph{interprétation de}~$P$ dans~$\mModel{M}$ si, pour tous
  préfixes~$\sigma,\tau\in P$, lorsque $\tau$ prolonge $\sigma$,
  on a $Rf(\sigma)f(\tau)$. En particulier, cela vaut lorsque
  $\tau=\sigma.n$.

  Relativement à une interprétation des préfixes~$P$, nous posons :
  \begin{enumerate}
  \item $\mModel{M}$ satisfait $\sFmla{\True}{!A}[\sigma]$ si et
    seulement si $\mSat{M}{!A}[f(\sigma)]$ ;
  \item $\mModel{M}$ satisfait $\sFmla{\False}{!A}[\sigma]$ si et
    seulement si $\mSat/{M}{!A}[f(\sigma)]$.
  \end{enumerate}
\end{defn}

La condition sur~$f$ s'applique aussi à toute extension utilisée
$\tau$ de~$\sigma$, souvent notée $\sigma.{*}$ : on a alors
$Rf(\sigma)f(\tau)$. Sur une branche construite à partir des
hypothèses de préfixe~$1$, les préfixes présents sont fermés par
préfixes initiaux ; une nouvelle extension~$\sigma.n$ ne possède
donc encore aucun descendant présent sur la branche.

\begin{defn}
  Soit $\Gamma$ un ensemble de !!{formula}s préfixées, et soit
  $P(\Gamma)$ l'ensemble des préfixes qui y figurent. Si $f$ est une
  interprétation de~$P(\Gamma)$ dans~$\mModel{M}$, nous disons que
  $\mModel{M}$ satisfait $\Gamma$ relativement à~$f$, et écrivons
  $\mSat{M}{\Gamma}[f]$, si $\mModel{M}$ satisfait chacune des
  !!{formula}s préfixées de~$\Gamma$ relativement à~$f$.
  L'ensemble~$\Gamma$ est \emph{satisfaisable} s'il existe un
  modèle~$\mModel{M}$ et une interprétation~$f$ de $P(\Gamma)$ tels
  que $\mSat{M}{\Gamma}[f]$.
\end{defn}

\begin{prop}
  Si $\Gamma$ contient $\sFmla{\True}{!A}[\sigma]$ et
  $\sFmla{\False}{!A}[\sigma.{*}]$ pour une !!{formula}~$!A$ et un
  préfixe~$\sigma$, ou contient
  $\sFmla{\True}{\lfalse}[\sigma]$, alors $\Gamma$ n'est pas
  satisfaisable.
\end{prop}

\begin{proof}
  Comme $\mSat/{M}{\lfalse}[f(\sigma)]$ pour tout monde, un
  ensemble~$\Gamma$ contenant
  $\sFmla{\True}{\lfalse}[\sigma]$ n'est pas satisfaisable.

  Supposons maintenant que $\Gamma$ contienne
  $\sFmla{\True}{!A}[\sigma]$ et
  $\sFmla{\False}{!A}[\tau]$, où $\tau$ est une extension utilisée
  de~$\sigma$. Si $\mSat{M}{!A}[f(\sigma)]$, alors, puisque
  $Rf(\sigma)f(\tau)$, la persistance de la vérité
  (\olref[sem][rel]{prop:true-monotonic}) donne
  $\mSat{M}{!A}[f(\tau)]$. Cela contredit
  $\mSat/{M}{!A}[f(\tau)]$.
\end{proof}

\begin{thm}[Correction]
  \ollabel{thm:tableau-soundness}
  Si $P(\Gamma)$ est fermé par préfixes initiaux non vides et si
  $\Gamma$ admet un !!{tableau} fermé, alors $\Gamma$ n'est pas
  satisfaisable.
\end{thm}

\begin{proof}
Une branche d'!!a{tableau} est \emph{satisfaisable} si et seulement
si l'ensemble des !!{signed formula}s qu'elle porte l'est ;
!!a{tableau} est satisfaisable s'il possède au moins une branche
satisfaisable.

Montrons que l'application d'une règle d'inférence à un !!{tableau}
satisfaisable produit encore un !!{tableau} satisfaisable. Cela
établira le théorème : tout !!{tableau} fermé résulte de règles
appliquées au !!{tableau} constitué des hypothèses de~$\Gamma$. Si
$\Gamma$ était satisfaisable, tout !!{tableau} ainsi obtenu le
serait également. Or un !!{tableau} fermé n'est pas satisfaisable,
car toutes ses branches sont fermées et donc non satisfaisables.

Supposons un !!{tableau} satisfaisable, c'est-à-dire !!a{tableau}
possédant au moins une branche satisfaisable.
L'application d'une règle ajoute des !!{signed formula}s à une
branche ou partage celle-ci en deux. Si le !!{tableau} possède une
branche satisfaisable que la règle ne touche pas, celle-ci reste
satisfaisable dans le nouveau
!!{tableau}. Il suffit donc d'examiner l'application d'une règle à
une branche satisfaisable.

Soit $\Gamma$ l'ensemble des !!{signed formula}s de cette branche,
et soit $\sFmla{S}{!A}[\sigma] \in \Gamma$ la !!{signed formula} à
laquelle la règle est appliquée. Sans partage, il faut montrer que
$\Gamma$ augmenté des conclusions reste satisfaisable. Avec
partage, il suffit de montrer qu'au moins l'une des deux nouvelles
branches est satisfaisable.
\begin{enumerate}
\item Appliquons $\TRule{\True}{\lnot}$ à
  $\sFmla{\True}{\lnot !B}[\sigma] \in \Gamma$. Pour toute
  extension utilisée~$\tau$ de~$\sigma$, la branche reçoit, parmi ses
  !!{signed formula}s,
  $\sFmla{\False}{!B}[\tau]$. Supposons
  $\mSat{M}{\Gamma}[f]$. Alors
  $\mSat{M}{\lnot !B}[f(\sigma)]$ ; aucun monde~$w$ tel que
  $Rf(\sigma)w$ ne force~$!B$. En particulier,
  $\mSat/{M}{!B}[f(\tau)]$. La même interprétation satisfait donc
  la branche prolongée.
\item Appliquons $\TRule{\False}{\lnot}$ à
  $\sFmla{\False}{\lnot !B}[\sigma] \in \Gamma$ : exercice.
\item Appliquons $\TRule{\True}{\land}$ à
  $\sFmla{\True}{!B \land !C}[\sigma] \in \Gamma$. La branche
  reçoit deux nouvelles !!{signed formula}s :
  $\sFmla{\True}{!B}[\sigma]$ et
  $\sFmla{\True}{!C}[\sigma]$. Si
  $\mSat{M}{\Gamma}[f]$, alors
  $\mSat{M}{!B \land !C}[f(\sigma)]$ ; donc
  $\mSat{M}{!B}[f(\sigma)]$ et
  $\mSat{M}{!C}[f(\sigma)]$. Le modèle satisfait les deux nouvelles
  formules signées relativement à~$f$.
\item Appliquons $\TRule{\False}{\lor}$ à
  $\sFmla{\False}{!B \lor !C}[\sigma] \in \Gamma$ : exercice.
\item Appliquons $\TRule{\False}{\lif}$ à
  $\sFmla{\False}{!B \lif !C}[\sigma] \in \Gamma$. La branche
  reçoit deux nouvelles !!{signed formula}s :
  $\sFmla{\True}{!B}[\sigma.n]$ et
  $\sFmla{\False}{!C}[\sigma.n]$, où $\sigma.n$ est nouveau sur la
  branche, c'est-à-dire $\sigma.n \notin P(\Gamma)$.

  Puisque $\Gamma$ est satisfaisable, il existe un modèle
  $\mModel{M}$ et une interprétation~$f$ de $P(\Gamma)$ tels que
  $\mSat{M}{\Gamma}[f]$, en particulier
  $\mSat/{M}{!B \lif !C}[f(\sigma)]$. La définition de
  l'implication fournit donc un monde~$w \in W$ tel que
  $Rf(\sigma)w$, $\mSat{M}{!B}[w]$ et
  $\mSat/{M}{!C}[w]$. Définissons $f'$ sur
  $P(\Gamma) \cup \{\sigma.n\}$ par $f'(\sigma.n)=w$ et
  $f'(\rho)=f(\rho)$ pour $\rho\in P(\Gamma)$. Les préfixes de la
  branche étant fermés par préfixes initiaux, aucun ancien préfixe ne
  prolonge $\sigma.n$ ; si $\rho$ est un préfixe initial de
  $\sigma.n$, la transitivité de~$R$ et $Rf(\sigma)w$ donnent
  $Rf'(\rho)f'(\sigma.n)$. Ainsi, $f'$ est une interprétation de
  $P(\Gamma) \cup \{\sigma.n\}$. Elle satisfait encore~$\Gamma$,
  et les deux nouvelles formules sont vraies au monde~$w$ avec leurs
  signes respectifs. Par conséquent,
  $\Gamma \cup \{\sFmla{\True}{!B}[\sigma.n],
  \sFmla{\False}{!C}[\sigma.n]\}$ est satisfaisable.
\end{enumerate}

Examinons maintenant les règles qui partagent une branche en deux.
\begin{enumerate}
\item Appliquons $\TRule{\False}{\land}$ à
  $\sFmla{\False}{!B \land !C}[\sigma] \in \Gamma$. La branche
  gauche reçoit $\sFmla{\False}{!B}[\sigma]$ et la droite
  $\sFmla{\False}{!C}[\sigma]$. Si
  $\mSat{M}{\Gamma}[f]$, alors
  $\mSat/{M}{!B \land !C}[f(\sigma)]$ ; il s'ensuit que
  $\mSat/{M}{!B}[f(\sigma)]$ ou
  $\mSat/{M}{!C}[f(\sigma)]$. Dans le premier cas, la branche
  gauche est satisfaisable ; dans le second, la droite l'est.
\item Appliquons $\TRule{\True}{\lor}$ à
  $\sFmla{\True}{!B \lor !C}[\sigma] \in \Gamma$ : exercice.
\item Appliquons $\TRule{\True}{\lif}$ à
  $\sFmla{\True}{!B \lif !C}[\sigma] \in \Gamma$ : exercice.
\end{enumerate}
\end{proof}

\begin{prob}
Complétez la preuve du
\olref[int][tab][sou]{thm:tableau-soundness}.
\end{prob}

\begin{cor}
\ollabel{cor:entailment-soundness}
Si $\Gamma \Proves !A$, alors $\Gamma \Entails !A$.
\end{cor}

\begin{proof}
  Si $\Gamma \Proves !A$, il existe $!B_1$, \dots, $!B_n \in
  \Gamma$ tels que
  $\Delta = \{\sFmla{\False}{!A}[1], \sFmla{\True}{!B_1}[1],
  \dots, \sFmla{\True}{!B_n}[1]\}$ admette un !!{tableau} fermé.
  Montrons que $\Gamma \Entails !A$. Supposons le contraire : il
  existe un modèle~$\mModel{M}$ et un monde~$w$ tels que
  $\mSat{M}{!B_i}[w]$ pour $i=1$, \dots,~$n$, mais
  $\mSat/{M}{!A}[w]$. Posons $f(1)=w$ ; alors $f$ interprète
  $P(\Delta)$ dans~$\mModel{M}$, et $\mModel{M}$ satisfait~$\Delta$
  relativement à~$f$. Mais, par le
  \olref{thm:tableau-soundness}, $\Delta$ n'est pas satisfaisable,
  puisqu'il admet un !!{tableau} fermé et que $P(\Delta)=\{1\}$ est
  fermé par préfixes initiaux non vides. Cette contradiction établit
  bien $\Gamma \Entails !A$.
\end{proof}

\begin{cor}
\ollabel{cor:weak-soundness}
Si $\Proves !A$, alors $!A$ est vraie dans tous les modèles.
\end{cor}

\end{document}
